% ============================================================
% Introduzione al formulario
% ============================================================
\section*{Come è fatto questo formulario}
\gdef\currentchapter{Introduzione}

Il formulario raccoglie le formule del corso, con il loro nome, e accompagna la dispensa. Ha due parti.
\begin{itemize}[itemsep=2pt]
  \item \textbf{Parte I, le formule.} Gli argomenti seguono l'ordine della dispensa, dal capitolo 1 al 26. Ogni voce
        ha un numero, un nome, la formula e, sotto, le \textbf{varianti} da usare nei casi particolari, ciascuna
        preceduta dal caso in cui vale: {\color{fcolore}\footnotesize$\blacktriangleright$}\ \textit{se $A$ e $B$
        sono incompatibili:}\ $P(A \cup B) = P(A) + P(B)$. In fondo, una \textbf{tabella} riassume le distribuzioni
        notevoli e i \textbf{richiami di matematica} raccolgono, per argomento, gli strumenti di analisi e di algebra
        lineare che servono.
  \item \textbf{Parte II, come si usano.} Per ogni formula c'è una scheda: \textbf{quando} si usa (le parole del
        testo che la fanno riconoscere), \textbf{come} si applica, quale \textbf{variante} scegliere, gli errori da
        evitare (\textbf{attenzione}) e un \textbf{esempio} con i numeri. Ogni argomento si apre con un riquadro che
        aiuta a scegliere la formula partendo dal testo dell'esercizio.
\end{itemize}

Il numero di ogni voce è un \textbf{collegamento}: nel PDF basta toccarlo per andare alla sua scheda d'uso, e il
titolo della scheda riporta alla formula. Per chi stampa, accanto c'è sempre il numero di pagina.

I colori sono quelli della dispensa: \textcolor{mspiform!80!black}{\textbf{verde}} per le formule,
\textcolor{mspiric!85!black}{\textbf{viola}} per i richiami di matematica, \textcolor{mspiteo!85!black}{\textbf{blu}}
per le schede d'uso, \textcolor{mspiB!85!black}{\textbf{arancione}} per gli errori tipici. L'etichetta \oltre{} segna
le poche formule che nelle slide non ci sono: generalizzazioni di quelle del corso, aggiunte perché negli esercizi
fanno comodo, e strumenti che la dispensa usa nelle dimostrazioni. La notazione è la stessa della dispensa:
$\compl{A}$ è il complementare di $A$, $\card{A}$ il numero dei suoi elementi, i vettori sono sottolineati una volta
($\vett{x}$) e le matrici due ($\matr{A}$).
